\documentclass[11pt]{article}
\usepackage[margin=1in]{geometry}
\usepackage{amsmath,amssymb,amsthm}
\usepackage{hyperref}
\usepackage{microtype}
\hypersetup{colorlinks=true,linkcolor=blue,citecolor=blue,urlcolor=blue}
\newtheorem{theorem}{Theorem}
\newtheorem{lemma}[theorem]{Lemma}
\newtheorem{proposition}[theorem]{Proposition}
\newtheorem{corollary}[theorem]{Corollary}
\theoremstyle{definition}
\newtheorem{definition}[theorem]{Definition}
\newcommand{\E}{\mathbb E}
\newcommand{\Prob}{\mathbb P}
\newcommand{\F}{\mathbb F}
\newcommand{\R}{\mathbb R}
\newcommand{\one}{\mathbf 1}
\DeclareMathOperator{\tr}{tr}
\DeclareMathOperator{\rank}{rank}
\DeclareMathOperator{\diag}{diag}
\DeclareMathOperator{\cum}{cum}
\newcommand{\norm}[1]{\left\lVert#1\right\rVert_2}
\title{A linear-width subspace embedding bound\\for the rerandomized SRHT}
\author{}
\date{}
\begin{document}
\maketitle

\begin{abstract}
For the normalized real Walsh transform with two independent sign
diagonals and uniform sampling without replacement, a universal
multiple of \(r/\varepsilon^2\) sampled columns suffices to embed every
fixed \(r\)-dimensional subspace with squared-norm distortion
\(1\pm\varepsilon\) and failure probability at most \(0.01\).
The proof bounds joint entry cumulants by connected graph contractions,
then bounds signed traces of a product of centered projections.
An explicit two-projection calculation converts these traces into
control of both spectral edges. All ranks and ambient Walsh dimensions
are included; the numerical constant is not optimized.
\end{abstract}

\section{Statement and notation}

Let \(n=2^m\), \(m\ge0\), and identify the coordinate set with
\(K=(\F_2)^m\). The real normalized Walsh matrix is
\[
 H_{ab}=n^{-1/2}(-1)^{\langle a,b\rangle},\qquad H^T=H,\quad H^2=I_n.
\]
Let \(D_1,D_2\) have mutually independent uniform Rademacher diagonal
entries. Independently let \(J\) be a uniformly chosen \(k\)-element
subset of \(K\), and let \(S_J\) collect its coordinate vectors. Put
\[
 \Omega_{n,k}=\sqrt{\frac nk}\,D_1HD_2HS_J.
\]
For a deterministic \(n\times r\) frame \(V\), \(V^TV=I_r\), write
\[
 X=HD_2HD_1V,\qquad P=XX^T,\qquad \delta=\frac rn.
\]
Thus \(X^TX=I_r\), and the Gram matrix in question is
\[
 V^T\Omega_{n,k}\Omega_{n,k}^TV
   =\frac nk X^T S_JS_J^T X.                                    \tag{1}
\]

\begin{theorem}\label{thm:main}
There is a universal constant \(C\ge1\) such that for every power of
two \(n\), every \(1\le r\le n\), and every \(0<\varepsilon<1\), one can
choose a deterministic integer
\[
 r\le k\le\min\{n,\lceil Cr/\varepsilon^2\rceil\}
\]
independently of \(V\), for which
\[
 \sup_{V^TV=I_r}
 \Prob\!\left\{\norm{V^T\Omega_{n,k}\Omega_{n,k}^TV-I_r}
                      >\varepsilon\right\}\le0.01.             \tag{2}
\]
For example, the proof permits
\[
 K_0=576\cdot4^{24},\qquad R_0=2^{20000},\qquad
 C=200R_0+8196K_0.                                               \tag{3}
\]
\end{theorem}

The supremum in (2) is outside the probability. The target frame is
fixed before both sign diagonals and the sample. No input permutation,
additional transform, or modified sampling law occurs in (1).
When \(k=n\), the Gram is exactly \(I_r\).
The printed assumption \(n=\Omega(\log r)\) in Problem~5.6 of
\cite{workshop} requires no additional restriction here, since
\(\log r\le r\le n\). Throughout, unadorned logarithms are natural,
and \(\log_2\) is the binary logarithm.

\section{Graph contractions and joint entry cumulants}

We use the following graph-operator result from Mingo and Speicher
\cite[Theorem~11 and Lemma~15]{mingo}.
A directed acyclic graph with prescribed source inputs and sink outputs,
every vertex on an input-to-output path, defines an operator by summing
its internal coordinate labels. Its norm is at most the product of the
edge-matrix norms. A connected graph without bridges can be modified
to have this form with any two prescribed distinct vertices as its
sole input and output. The modifications transpose reversed edges and
split vertices using identity edges; they preserve the graph sum and
the product of edge norms. Vertex-dependent dimensions and rectangular
edge matrices are allowed.

The following rank refinement makes the use of that result explicit.
A loop counts twice toward a vertex degree.

\begin{lemma}\label{lem:graph}
Let a connected finite multigraph have positive even degrees. Assign
a matrix of operator norm at most one to each edge, with compatible
coordinate dimensions at its endpoints. Suppose one edge matrix is a
real rank-\(r\) orthogonal projection \(P_0=VV^T\), \(r\ge1\).
At each vertex also allow a coordinate weight of absolute value at most
one. The unrestricted signed sum of the products of all these entries
and weights over the vertex labels has absolute value at most \(r\).
\end{lemma}
\begin{proof}
A connected even-degree graph has no bridge: a one-edge cut would give
an odd degree sum on one side. Represent each vertex weight as a
diagonal loop of norm at most one. Replace the distinguished edge
entry by
\[
 (P_0)_{ij}=\sum_{a=1}^r V_{ia}V_{ja}.
\]
This inserts a degree-two vertex of dimension \(r\). Split that vertex
into two dimension-\(r\) vertices, put one incident edge at each, and
join them by \(I_r\). The signed sum is unchanged. The resulting graph
is still connected with even degrees, also when the original edge was
a loop.

Take the two new vertices as the input and output in the cited
graph-operator result. The resulting operator \(T:\R^r\to\R^r\) has
norm at most one, since both rectangular factors and every inserted
identity have norm one. The full signed sum is
\(\one_r^T T\one_r\), of absolute value at most \(r\).
This uses actual dimension-\(r\) vertices; it does not replace an
ambient dimension by rank without a factorization.
\end{proof}

For scalar variables with the requisite moments, define
\[
 \cum(Y_1,\ldots,Y_q)=
 \sum_{\pi\in\mathcal P([q])}(-1)^{|\pi|-1}(|\pi|-1)!
                       \prod_{B\in\pi}\E\prod_{j\in B}Y_j.       \tag{4}
\]
The moment-cumulant formula is the inverse identity.
We will use both only as finite polynomial identities.

\begin{lemma}\label{lem:cumulants}
For the projection \(P=HD_2HD_1VV^TD_1HD_2H\), arbitrary deterministic
row pairs \((i_f,j_f)\), and every integer \(q\ge1\),
\[
 \left|\cum(nP_{i_1j_1},\ldots,nP_{i_qj_q})\right|
                      \le (2q)^{12q}r.                         \tag{5}
\]
The cumulant is zero unless
\[
                 \sum_{f=1}^q(i_f+j_f)=0\quad\hbox{in }K.        \tag{6}
\]
Moreover, \(\E P=\delta I_n\).
\end{lemma}
\begin{proof}
We first recall the finite cumulant formula for products. If each
\(Y_f\) is a product of underlying variable occurrences, let \(\tau\)
group the occurrences into those products. Expand (4) by the
moment-cumulant formula for the underlying variables. A partition
\(\sigma\) of their occurrences survives exactly when
\(\sigma\vee\tau\) has one block, with coefficient one. Indeed, if its
induced partition of the \(q\) products has \(b\) blocks, the remaining
coefficient is
\[
 \sum_{j=1}^b S(b,j)(-1)^{j-1}(j-1)!=1_{\{b=1\}},
\]
where \(S(b,j)\) is the number of partitions into \(j\) blocks.
This identity follows by comparing coefficients in
\(\log(\exp z)=z\). Joint cumulants crossing independent families
vanish, as follows by taking the logarithm of the factorized formal
moment series.

Write \(P_0=VV^T\) and \(\chi_i(a)=(-1)^{\langle i,a\rangle}\).
The exact entry expansion is
\[
 nP_{ij}=\sum_{a,b,x,y}
 \chi_i(a)\chi_j(b)\,
 \xi_{2,a}\xi_{2,b}\xi_{1,x}\xi_{1,y}
                  H_{ax}(P_0)_{xy}H_{yb}.                       \tag{7}
\]
The factor \(n\) cancels the two outer Hadamard normalizations.
Both inner Hadamard matrices in (7) remain normalized.

Apply the product formula to \(q\) copies of (7). A nonzero partition
has a partition \(\alpha\) of its \(2q\) middle-sign occurrences and
a partition \(\beta\) of its \(2q\) input-sign occurrences, all blocks
even. Equal-coordinate occurrences within a block are identified.
Different blocks may receive equal labels; no inequality between
their labels is imposed. The connected-join condition connects all
\(q\) matrix paths \(H\)--\(P_0\)--\(H\).

The resulting graph has one vertex per middle-sign block, of degree
equal to its even size, and one per input-sign block, of degree twice
its size. It is connected with positive even degrees. Its matrices
are normalized \(H\)'s and \(P_0\)'s, all of norm one. The outer
characters become vertex weights of modulus one.
Lemma~\ref{lem:graph} bounds its unrestricted signed label sum by \(r\).
Absolute values have not been inserted into the matrix-entry products.

The order-\(b\) cumulant of a single Rademacher sign has absolute value
at most \(b^{2b}\): in (4) there are at most \(b^b\) partitions, each
factorial coefficient is at most \(b!\le b^b\), and each moment has
absolute value at most one. There are at most \((2q)^{4q}\) pairs of
layer partitions. On each layer, the product of cumulant coefficients
is at most \((2q)^{4q}\), since the block sizes sum to \(2q\).
Both layers and their count prove (5).

For (6), let \(C_s=\diag(\chi_s(i))\) and let \(T_s\) translate
coordinates by \(s\). The Walsh identities give
\[
 C_sH=HT_s,\qquad T_sH=HC_s,\qquad
 C_sHD_2HD_1V=H(T_sD_2T_s)H(C_sD_1)V.                           \tag{8}
\]
The two new diagonals remain independent iid sign diagonals. Thus
\(C_sPC_s\) has the same law as \(P\) for the same fixed \(V\).
Multilinearity multiplies the indicated cumulant by
\(\chi_s(\sum_f(i_f+j_f))\). Character separation proves (6).

Finally set \(B=HD_1P_0D_1H\). Averaging \(D_2\) retains only
\(\diag(B)\) between the outer \(H\)'s, and averaging \(D_1\) makes
each of those diagonal entries \(r/n\). Hence \(\E P=\delta I_n\).
\end{proof}

Set \(R=P-\delta I_n\). Every entry of \(R\), including every diagonal
entry, has mean zero. Cumulants of order at least two are unchanged
by this centering, so (5)--(6) hold for \(R\) at those orders.
The zero diagonal mean will be used in the next proof.

\section{A signed trace bound for centered projections}

\begin{proposition}\label{prop:trace}
Let \(E\) be a diagonal projection with independent Bernoulli entries
of mean \(\theta\), independent of \(P\). Put
\[
 W=E-\theta I_n,\quad A=RW,\quad \kappa=n\theta,\quad
 K_2=4^{24},\quad K_0=576K_2.
\]
For an integer \(p\ge2\), suppose
\[
 0<\theta\le\tfrac12,\qquad \kappa\ge r,\qquad
                         r\ge2(4p+1)^{1000}.                    \tag{9}
\]
Then
\[
                  \left|\E\tr A^{2p}\right|
                           \le K_0^p(\delta\theta)^p.            \tag{10}
\]
\end{proposition}

The absolute value in (10) is outside the expectation. In particular,
(10) is not a norm moment of the generally nonsymmetric matrix \(A\).

\subsection{Selector equality graphs}

Expand the cyclic trace and first average the independent selectors.
A partition of its \(2p\) cyclic positions specifies the equal row
labels. Singleton blocks vanish. For a surviving partition, write
\(m_u\ge2\) for its block sizes and
\[
 v=p-s,\quad 0\le s\le p-1.
\]
Its quotient multigraph is connected, with \(2p\) edge occurrences
and degrees \(2m_u\ge4\), so
\[
                         \sum_u(\deg u-4)=4s.                   \tag{11}
\]
The graph contraction sums over distinct full row labels for its
\(v\) vertices. Let \(t\) count the vertices incident to a
\emph{nonloop} unordered edge of odd multiplicity, and let \(\ell\)
count loop occurrences. Thus \(s+t\le p\).
For a centered Bernoulli variable \(w\) and \(b\ge2\),
\[
 |\E w^b|
 \le\theta(1-\theta)\big((1-\theta)^{b-1}+\theta^{b-1}\big)
 \le\theta.                                                     \tag{12}
\]
The absolute selector coefficient of a graph is therefore at most
\(\theta^v\). These are ordinary moments of equality blocks, not
selector cumulants.

\begin{lemma}\label{lem:count}
Let \(N_{p,s,t}\) count these equality partitions, and put \(L=4p+1\).
Then
\[
 \ell\le4t+12s+2,\qquad
                N_{p,s,t}\le8p\,3^p L^{74(s+t)}.                 \tag{13}
\]
\end{lemma}
\begin{proof}
Mark all odd-incident vertices and all vertices of degree greater
than four, and call their union \(B\). There are at most \(2s\)
vertices of the latter kind. By (11),
\[
 |B|\le t+2s,\qquad D_B:=\sum_{u\in B}\deg u\le4t+12s.            \tag{14}
\]
Outside \(B\), degree is four and all nonloop multiplicities are
even. In a graph with more than one vertex, the possible types are:
two distinct neighbors joined by doubled edges and no loop;
one neighbor joined by a doubled edge and one loop; or one neighbor
joined by four edges and no loop. Call the first type internal and
the last two terminal.

When \(B\ne\varnothing\), a four-edge terminal is directly adjacent
to \(B\); otherwise it would form an isolated two-vertex component.
Every loop terminal joins \(B\) by a doubled path of internal
vertices. Different terminal attachments consume distinct pairs
of half-edges at \(B\). Loops in \(B\) themselves consume two
half-edges each. In particular the total number of loops is at most
\(D_B\). When \(B=\varnothing\), the graph is a doubled cycle
(including four parallel edges on two vertices) or a doubled path
with a loop at each endpoint. The one-vertex possibility has two
loops. This proves the loop bound in (13).

For the count, first suppose \(a=s+t\ge1\), so \(B\ne\varnothing\).
Retain \(B\) and all terminal vertices, and contract maximal chains
of internal doubled-path vertices into distinguished doubled links.
A link may return to the same retained vertex or be parallel to
another edge; its paired half-edge ports must be retained.
Direct doubled edges may be given length zero. Connectivity excludes
an isolated cycle of removed vertices.

If \(T\) terminal vertices are retained, then \(T\le D_B/2\).
The retained graph has \(q\) vertices and total degree \(D\), with
\[
 q\le |B|+D_B/2\le8a,\qquad
 D=D_B+4T\le3D_B\le36a,\qquad q\le p,\quad D\le4p.                \tag{15}
\]
Temporarily label its \(q\) vertices and distinguish their half-edges.
The following data give a surjective encoding of its possible
expansions.

First choose \(q\) and its degree vector. This costs at most
\(L^aL^{8a}\). Pair its \(D\) half-edges to specify the retained
multigraph, allowing loops, in at most \(D^{D/2}\le L^{18a}\)
ways. At each vertex designate disjoint pairs of half-edges as
doubled-link ports. Giving each half-edge either no mate or a named
mate costs at most \((D+1)^D\le L^{36a}\); retain only compatible
choices. The edge pairing then identifies each doubled link through
its ports, even when it is parallel to an odd edge or returns to one
vertex through two different ports. There are at most \(D/4\le9a\)
links. Assign each a length between zero and \(p\), costing at most
\(L^{9a}\), and expand its internal vertices in path order.
Keep only expansions with the prescribed vertex and edge counts.
These choices give at most \(L^{72a}\) graph encodings.

Every original graph yields such data by its specified contractions.
Conversely, valid data reconstruct its graph with distinguished
half-edges. Invalid and duplicate codes enlarge the count. Core
labels are auxiliary; no separately labeled set of original vertices
is chosen.

An Euler circuit pairs arriving and departing half-edges at each
vertex. For degree \(2m_u\), the number of transitions is bounded by
\[
 (2m_u-1)!!=3\prod_{j=3}^{m_u}(2j-1)
                                  \le3(4p)^{m_u-2}.
\]
Their product is at most \(3^p L^{2s}\), because
\(\sum_u(m_u-2)=2s\). Choose a starting directed edge in at most
\(4p\) ways. A transition system producing one circuit determines
the cyclic word, hence its equality partition after vertex names
are discarded. Every original partition is produced. No additional
choice of the positions of retained vertices is needed.
The result is at most \(4p\,3^p L^{74a}\).

For \(a=0\), there are at most the two underlying graphs described
above. All degrees are four, so the same circuit encoding costs
at most \(8p\,3^p\). This proves (13).
\end{proof}

\subsection{Entry-cumulant partitions and their counts}

Fix one equality graph. Expand the expectation of its \(2p\) entries
of \(R\) by entry cumulants. Singleton blocks vanish, including
singleton loops. Let a surviving edge partition have
\[
                         b=p-d,\qquad 0\le d\le p-1
\]
blocks. Its binary constraint matrix has one row per block: the
sum of that block's endpoint incidence vectors over \(\F_2\).
Write \(h\) for its rank. By Lemma~\ref{lem:cumulants}, a summand
can be nonzero only on its XOR constraints, which have
\(n^{v-h}\) unrestricted solutions.

Two kinds of pair block vanish on every injective labeling:
adjacent nonparallel nonloop edges force equality of their other
endpoints, and a loop paired with a nonloop edge forces equality
of that edge's endpoints. Delete these partitions before relaxing
injectivity. A remaining pair consists of parallel nonloop
occurrences, disjoint nonloop edges, or two loops. Only a disjoint
nonloop pair has a nonzero binary row, of weight four.

Let \(M\) count occurrences in blocks of size at least three.
Their total excess over size two is \(2d\), whence
\[
 M\le6d.
\]
At most \(12d\) vertices touch such occurrences. Any other odd-incident
vertex must occur in a disjoint-pair row: an odd number of its chosen
edge's occurrences cannot all pair among themselves, and the forbidden
pair types have been removed. The supports of the disjoint-pair rows
lie in the union of the supports of a basis chosen from those
weight-four rows. Their rank is at most \(h\). Consequently
\[
                              t\le12d+4h.                       \tag{16}
\]

We count the possible entry partitions with fixed \(d,h\), keeping
the degree-excess budget. Choose the at most \(6d\) large-block
occurrences and partition them in at most
\[
                 (2p+1)^{6d}(2p)^{6d}\le L^{12d}
\]
ways. Their cumulant constants from (5) are at most
\((4p)^{12M}\le L^{72d}\). Pair blocks use \(K_2=4^{24}\), so
their constants together are at most \(K_2^p\).

Let \(U\) be the union of the supports of disjoint-pair rows.
It has at most \(4h\) vertices and can be chosen in at most
\(L^{4h}\) ways. Every occurrence in such a pair has both endpoints
in \(U\). The number available is at most
\[
 \frac12\sum_{u\in U}\deg u\le2|U|+2s\le8h+2s.
\]
Choosing those occurrences and pairing them costs at most
\[
 2^{8h+2s}(2p)^{4h+s}\le L^{8h+2s}.
\]
Including \(U\), this is \(L^{12h+2s}\).

All remaining loop occurrences pair with each other, at a cost
bounded by
\[
                  (2p)^{\ell/2}\le L^{2t+6s+1}.                 \tag{17}
\]
For a nonloop unordered edge with positive even remaining
multiplicity \(a_e\), its parallel-pair choices are at most
\(3(4p)^{(a_e-4)_+/2}\). There are at most \(p\) such groups.
At each vertex, the sum of \((a_e-4)_+\) over incident groups is
at most \(\deg u-4\): if there is any such group, its total
multiplicity is bounded by that degree, and at least one four
has been subtracted. Summing over endpoints and using (11) gives
\(\sum_e(a_e-4)_+\le2s\). All parallel-pair choices therefore
cost at most \(3^pL^s\).

The total partition count and cumulant constants for fixed \(d,h\)
are thus bounded by
\[
                         (3K_2)^p L^{84d+12h+9s+2t+1}.           \tag{18}
\]
These counts permit incompatible choices but cover every admissible
partition. No common sign is assumed for cumulant terms.

Each cumulant block contributes one factor \(r\), and the entry
normalizations contribute \(n^{-2p}\). After deleting the
pointwise-vanishing partitions, relax injectivity in the absolute
upper bound and use \(n^{v-h}\) labels. The selector coefficient
from (12) is at most \(\theta^v\). Their combined dimension factor
is exactly
\[
 \theta^v r^{p-d}n^{-2p}n^{v-h}
            =(\delta\theta)^p\kappa^{-s}r^{-d}n^{-h}.             \tag{19}
\]
The factor \(\kappa^{-s}\) is essential and has not yet been discarded.

\subsection{Summation and proof of Proposition~\ref{prop:trace}}

Multiply (18)--(19) by (13). Finite expansion followed by the stated
triangle inequalities bounds \(|\E\tr A^{2p}|\) by
\[
 8pL(9K_2)^p(\delta\theta)^p
  \sum_{s,t,d,h}L^{83s+76t+84d+12h}\kappa^{-s}r^{-d}n^{-h},       \tag{20}
\]
where (16) holds. For fixed \(s,d,h\),
\[
 \sum_{t=0}^{12d+4h}L^{76t}\le2L^{912d+304h}.
\]
Use \(\kappa\ge r\), \(n\ge r\), and extend the other parameter
ranges to nonnegative integers. Since the resulting exponents
are \(83s+996d+316h\), (20) is at most
\[
 16pL(9K_2)^p(\delta\theta)^p
          \sum_{s,d,h\ge0}\left(\frac{L^{1000}}r\right)^{s+d+h}.
\]
Under (9) the three geometric series have product at most eight.
Finally
\[
 128pL\le640p^2\le64^p\qquad(p\ge2);
\]
the last inequality holds at \(p=2\), and its successive left-hand
ratio is smaller than \(64\). This proves (10).

\section{From signed traces to both spectral edges}

\begin{lemma}\label{lem:transfer}
Let \(P=XX^T\) and \(E\) be any real orthogonal projections,
\(X^TX=I_r\). Let \(0<\delta,\theta\le1/2\) and set
\[
 A=(P-\delta I)(E-\theta I),\qquad
 a^2=\delta(1-\delta)\theta(1-\theta).
\]
For every integer \(p\ge1\), \(\tr A^{2p}\ge-2r a^{2p}\).
If \(0<\eta<1\) and \(\delta\le\theta\eta^2/64\), then pointwise
\[
 1_{\{\norm{\theta^{-1}X^TEX-I_r}>\eta\}}
          (\theta\eta/2)^{2p}
                \le\tr A^{2p}+2r a^{2p}.                        \tag{21}
\]
\end{lemma}
\begin{proof}
Diagonalize \(PEP\) on \(\operatorname{range}P\).
For a unit eigenvector \(u\) with eigenvalue \(0<\lambda<1\), put
\[
 w=\frac{Eu-\lambda u}{\sqrt{\lambda(1-\lambda)}}.
\]
Then \(Pw=0\), \(\norm w=1\), and the matrices on the orthonormal
basis \(u,w\) are
\[
 P=\begin{pmatrix}1&0\\0&0\end{pmatrix},\qquad
 E=\begin{pmatrix}\lambda&\sqrt{\lambda(1-\lambda)}\\
                  \sqrt{\lambda(1-\lambda)}&1-\lambda\end{pmatrix}.
                                                                  \tag{22}
\]
Indeed \(PEu=\lambda u\), \(\|Eu\|^2=\lambda\), and \(E^2=E\)
give these assertions directly. Distinct eigenvectors from an
orthonormal eigenbasis give orthogonal pairs, since the relevant
cross inner products reduce to \(\langle u,Eu'\rangle=0\),
also in repeated eigenspaces. The displayed spaces are invariant.

Eigenvalues \(\lambda=0,1\) give the intersections of
\(\operatorname{range}P\) with \(\ker E,\operatorname{range}E\),
respectively. The remaining orthogonal complement lies in \(\ker P\),
is invariant under \(E\), and can be diagonalized there.
This constructs the entire decomposition directly.

On (22), the characteristic polynomial of \(A\) is
\[
 z^2-tz+a^2,\qquad t=\lambda-\delta-\theta+2\delta\theta.           \tag{23}
\]
If its roots are nonreal, they are conjugates of modulus \(a\),
and their even-power sum is at least \(-2a^{2p}\).
If they are real, that sum is nonnegative.
Power traces equal the sums of root powers even at a defective
repeated root, by the Cayley--Hamilton recurrence.
The remaining one-dimensional eigenvalues are among
\[
 -(1-\delta)\theta,\quad (1-\delta)(1-\theta),\quad
                  -\delta(1-\theta),\quad \delta\theta,
\]
whose even powers are nonnegative. There are at most \(r\)
two-dimensional blocks. This proves the lower trace bound.

On the bad event in (21), some compressed eigenvalue satisfies
\(|\lambda-\theta|>\theta\eta\). For an interior \(\lambda\), set
\(b=\theta\eta\). The stated hypothesis gives
\[
 |t|>63b/64,\qquad a\le\sqrt{\delta\theta}\le b/8.
\]
Both roots in (23) are therefore real and have the same sign.
Their smaller absolute value is at most \(a\), so the larger one
exceeds \(|t|-a>55b/64>b/2\). This block contributes more than
\((b/2)^{2p}\), while all other blocks together contribute at
least \(-2r a^{2p}\).
For \(\lambda=0\), the real eigenvalue has magnitude
\((1-\delta)\theta\ge\theta/2>b/2\).
For \(\lambda=1\), its magnitude is at least \(1/4>b/2\).
These include \(E=0\) and \(E=I\).
On the complementary event the lower trace bound makes the right
side of (21) nonnegative, proving the pointwise inequality.
\end{proof}

\begin{corollary}\label{cor:bernoulli}
Suppose
\[
 p=\lceil\log_2(3000r)\rceil,\qquad r\ge2(4p+1)^{1000}.
\]
For \(0<\eta<1\), \(0<\theta\le1/2\), and
\(\kappa=n\theta\ge64K_0r/\eta^2\), the exact random projection \(P\)
and independent Bernoulli coordinate projection \(E\) satisfy
\[
 \Prob\{\norm{\theta^{-1}X^TEX-I_r}>\eta\}\le1/1000.              \tag{24}
\]
\end{corollary}
\begin{proof}
The width hypothesis implies \(\kappa\ge r\) and
\(\delta\le\theta\eta^2/64\), so both preceding results apply.
Taking expectations in (21) and using \(a^2\le\delta\theta\) gives
\[
 \Prob(\text{bad})
 \le(K_0^p+2r)\left(\frac{4r}{\kappa\eta^2}\right)^p
 \le(1+2r)16^{-p}\le3r\,2^{-p}\le1/1000.
\]
No positivity of the unshifted signed trace was assumed.
\end{proof}

\section{Fixed-size sampling and the remaining dimensions}

\begin{lemma}\label{lem:sampling}
Let \(X\) be an orthonormal frame, possibly random and independent
of sampling. For \(0<\varepsilon<1\), an integer \(1\le k\le n\),
and
\[
 \alpha=\varepsilon/4,\qquad
 \theta_\pm=(1\pm\alpha)k/n\le1,
\]
suppose each normalized Bernoulli Gram of density \(\theta_\pm\)
has error at most \(\varepsilon/4\) except with probability
\(\gamma\). A uniform \(k\)-element sample then satisfies
\[
 \Prob\left\{\norm{\frac nk X^TS_JS_J^TX-I_r}>\varepsilon\right\}
 \le2\gamma+2\exp(-\varepsilon^2 k/48).                           \tag{25}
\]
\end{lemma}
\begin{proof}
Take independent uniform \(U_i\) on \([0,1]\), independent of \(X\).
Let \(J\) index the \(k\) smallest values, and put
\(B_\pm=\{i:U_i\le\theta_\pm\}\), \(N_\pm=|B_\pm|\).
Their indicators have the required Bernoulli marginal laws.
On \(N_-\le k\le N_+\), set inclusion and positivity give
\[
 (1-\alpha)\theta_-^{-1}X^TE_{B_-}X
 \preceq \frac nkX^TE_JX
 \preceq (1+\alpha)\theta_+^{-1}X^TE_{B_+}X.
\]
On the two Gram success events its spectrum lies between
\((1-\alpha)^2\) and \((1+\alpha)^2\), which are within
\(1\pm\varepsilon\).

For a binomial \(N\) of mean \(\mu\), its exponential moment gives
\[
 \Prob(N-\mu\ge a)\le e^{-a^2/(2\mu+2a/3)},\qquad
 \Prob(N-\mu\le-a)\le e^{-a^2/(2\mu)}.                            \tag{26}
\]
For completeness,
\(\E e^{t(N-\mu)}\le\exp(\mu(e^t-1-t))\).
The bound \(e^t-1-t\le t^2/(2(1-t/3))\), \(0<t<3\),
follows termwise from \(j!\ge2\cdot3^{j-2}\), \(j\ge2\).
Choose \(t=a/(\mu+a/3)\) for the upper tail.
For the lower tail use \(e^{-t}-1+t\le t^2/2\) and \(t=a/\mu\).
This proves (26).

For \(N_-\), use \(\mu=(1-\alpha)k\), \(a=\alpha k\), giving
an exponent at least \(\alpha^2k/2\). For \(N_+\), the lower-tail
exponent is at least \(\alpha^2k/3\), since \(\alpha<1/4\).
The two count failures sum to at most
\(2e^{-\alpha^2k/3}\). Union bounding proves (25).
This argument does not assume monotonicity of normalized failure
probabilities as a function of \(k\).
\end{proof}

\begin{lemma}\label{lem:small}
For every deterministic orthonormal frame \(V\), let \(X=HD_2HD_1V\)
and let \(J\) be an independent uniform \(k\)-element subset. Then
\[
 \E_{D_1,D_2,J}\left\|\frac nkX^TS_JS_J^TX-I_r\right\|_F^2
                                      \le\frac{r(r+1)}k.         \tag{27}
\]
Consequently \(k=\min\{n,\lceil200r^2/\varepsilon^2\rceil\}\)
has failure probability at most \(0.01\).
\end{lemma}
\begin{proof}
The case \(n=1\), and every case \(k=n\), is exact.
For \(n>1\), let \(x_i^T\) be the rows of \(X\).
The one- and two-coordinate inclusion probabilities of a uniform
sample, together with \(\sum_i x_ix_i^T=I_r\), give
\[
 \E\!\left[\|G-I_r\|_F^2\mid X\right]
  =\frac{n-k}{k(n-1)}
                \left(n\sum_i\|x_i\|^4-r\right).                \tag{28}
\]
This follows by expanding the ordered products, using
\(\sum_{i\ne j}(x_ix_i^T)(x_jx_j^T)
=I_r-\sum_i\|x_i\|^2x_ix_i^T\), and taking the trace.

Condition on \(U=HD_1V\), with rows \(u_j^T\).
Every \(\sqrt n\,x_i\) is a Rademacher sum of these rows.
The three fourth-moment pairings, subtracting twice the all-equal
pattern, give
\[
 \E_{\xi}\left\|\sum_j\xi_j u_j\right\|^4
             =r^2+2r-2\sum_j\|u_j\|^4\le r^2+2r.
\]
Average (28) over \(D_2\) using this bound, and then over \(D_1\).
Since \((n-k)/(n-1)\le1\), this gives (27).
Markov's inequality and the domination of operator norm by
Frobenius norm give failure at most
\(r(r+1)/(k\varepsilon^2)\).
At the stated nonfull width this is at most
\((r+1)/(200r)\le0.01\); full width is exact.
\end{proof}

\subsection{Proof of Theorem~\ref{thm:main}}

We use the constants in (3). First, for every \(r\ge R_0\), the
integer \(p=\lceil\log_2(3000r)\rceil\) satisfies the rank condition
of Corollary~\ref{cor:bernoulli}. To check this explicitly, set
\(x=\log_2r\ge20000\). Then \(p\le x+13\), so \(4p+1\le4x+53\).
The function
\[
 f(x)=x-1-1000\log_2(4x+53)
\]
is positive at \(x=20000\), since \(80053<2^{17}\), and is
increasing thereon:
\(f'(x)=1-4000/((4x+53)\log 2)>0\).
For the last inequality, \(\log2>1/2\) already suffices.
Thus \(r\ge2(4p+1)^{1000}\).

If \(r<R_0\), choose the width in Lemma~\ref{lem:small}.
It obeys the desired probability and
\[
 k\le\min\{n,\lceil200R_0r/\varepsilon^2\rceil\}.
\]

For \(r\ge R_0\), put
\[
                    k_0=\left\lceil
                            \frac{2048K_0r}{\varepsilon^2}
                          \right\rceil.                         \tag{29}
\]
If \(n\le4k_0\), choose \(k=n\). Its Gram is exact, and
\[
 k=n\le4k_0\le(8192K_0+4)r/\varepsilon^2
                         \le8196K_0r/\varepsilon^2.             \tag{30}
\]
Here \(r/\varepsilon^2\ge1\) and \(K_0\ge1\).

Otherwise choose \(k=k_0\), and take
\(\theta_\pm=(1\pm\varepsilon/4)k_0/n\).
Because \(n>4k_0\), both are at most \(5/16<1/2\).
Their effective widths satisfy
\[
 n\theta_\pm\ge\frac34 k_0
 \ge\frac{1536K_0r}{\varepsilon^2}
 \ge\frac{64K_0r}{(\varepsilon/4)^2}.                             \tag{31}
\]
Apply Corollary~\ref{cor:bernoulli} with \(\eta=\varepsilon/4\).
Both Bernoulli failure probabilities are at most \(1/1000\).
Lemma~\ref{lem:sampling} bounds the fixed-size failure by
\[
             \frac2{1000}+2e^{-\varepsilon^2k_0/48}<0.01.        \tag{32}
\]
Indeed \(\varepsilon^2k_0/48\ge2048K_0r/48>42\), and
\(e^{42}\ge1+42+42^2/2=925>250\), which suffices for (32).

All choices are deterministic functions of \(n,r,\varepsilon\).
They satisfy \(k\ge r\), \(k\le n\), and the upper width bound with
the single constant \(C\) in (3). Every preceding probability estimate
is uniform in the deterministic orthonormal frame \(V\).
Equation (1) therefore gives the supremum bound (2), completing
the positive resolution. \qed

\begin{thebibliography}{9}
\bibitem{mingo}
J.~A. Mingo and R.~Speicher,
\emph{Sharp Bounds for Sums Associated to Graphs of Matrices},
arXiv:0909.4277 (2009).
\url{https://arxiv.org/abs/0909.4277}.

\bibitem{workshop}
Amsel et al.,
\emph{Linear Systems and Eigenvalue Problems:
Open Questions from a Simons Workshop},
arXiv:2602.05394v3 (2026), Definitions~5.1 and~5.3 and Problem~5.6.
\url{https://arxiv.org/abs/2602.05394v3}.
\end{thebibliography}
\end{document}
